\section{Design}

GlassVM defines a shared orchestration boundary around machine-owned
semantics. The boundary is expressed as a small set of contracts rather than
as a universal emulator implementation.

\paragraph{Machine identity and admission.}
\texttt{MachineDescriptor} identifies a bundle and describes the execution
coordinates, artifact expectations, and supported machine semantics exposed
to callers. \texttt{MachineContract} supplies machine-specific admission
rules, configuration and input schemas, and the observation capabilities the
bundle can provide. These declarations let orchestration validate a request
without transferring machine semantics into the facade.

\paragraph{Execution request and preparation.}
An \texttt{ExecutionRequest} keeps artifact bytes, typed configuration,
scheduled inputs, execution controls, and the observation request as separate
values. \texttt{EmulatorBackend} validates the execution-side values and
constructs a session only after preparation succeeds. Preparation returns
two distinct objects: \texttt{PreparedRun} is validated execution-side state
and identity; \texttt{PreparedObservation} is the prospective, negotiated
evidence contract. The pair makes it possible to reason separately about
which machine run is admitted and which observations that run will attempt to
produce.

\paragraph{Session and evidence derivation.}
An \texttt{EmulatorSession} advances the selected machine and emits typed
records at supported boundaries. A \texttt{Normalizer} declares capability
dependencies and reduces selected evidence as it arrives. It receives a
begin notification, zero or more observations, and finalization with any
negotiated final state. The normalizer owns only its reducer state; the
framework does not create a hidden evidence bag or require a whole-run trace.

\paragraph{Delivery and accounting.}
An \texttt{EmissionSink} is caller-owned and consumes each synchronous batch.
The \texttt{Evidence\allowbreak Receipt} retrospectively describes fulfillment
of the negotiated observation request. A \texttt{Recorder\allowbreak Receipt}
describes what the reference recorder persisted and whether its package can
be published. A \texttt{FileRun\allowbreak Session} coordinates execution,
sink delivery, recorder finalization, and publication while preserving these
independent outcomes.

The lifecycle and ownership boundaries are summarized in
Figure~\ref{fig:architecture}. The same emission stream can be delivered to a
Python consumer and to the Rust file recorder; these consumers do not define
separate machine execution paths.

\begin{figure}[t]
  \centering
  \resizebox{0.90\linewidth}{!}{%
\begin{tikzpicture}[
  >=Stealth,
  box/.style={draw=black!70, rounded corners, fill=gray!8, align=center,
              text width=3.0cm, minimum height=0.65cm, inner sep=4pt},
  flow/.style={->, thick, draw=black!75}
]
  \node[box] (experiment) at (0, 0) {Experiment};
  \node[box] (request) at (0, -1.0) {ExecutionRequest + ObservationRequest};
  \node[box] (prepare) at (0, -2.0) {prepare};
  \node[box] (run) at (-2.1, -3.0) {PreparedRun};
  \node[box] (observation) at (2.1, -3.0) {Prepared\\Observation};
  \node[box] (session) at (0, -4.0) {EmulatorSession};
  \node[box] (emissions) at (0, -5.0) {Bounded emission stream};
  \node[box] (python) at (-2.1, -6.0) {Python consumers};
  \node[box] (recorder) at (2.1, -6.0) {Rust recorder};
  \node[box] (downstream) at (-2.1, -7.0) {Downstream analysis};
  \node[box] (receipt) at (2.1, -7.0) {RecorderReceipt};
  \node[box] (published) at (2.1, -8.0) {PublishedRun};

  \draw[flow] (experiment) -- (request);
  \draw[flow] (request) -- (prepare);
  \draw[flow] (prepare) -- (run);
  \draw[flow] (prepare) -- (observation);
  \draw[flow] (run) -- (session);
  \draw[flow] (observation) -- (session);
  \draw[flow] (session) -- (emissions);
  \draw[flow] (emissions) -- (python);
  \draw[flow] (emissions) -- (recorder);
  \draw[flow] (python) -- (downstream);
  \draw[flow] (recorder) -- (receipt);
  \draw[flow] (receipt) -- (published);
\end{tikzpicture}%
}

  \caption{Preparation separates validated execution state from negotiated
  observations. A synchronous emission stream is delivered to caller-owned
  consumers and, when requested, the reference Rust recorder.}
  \label{fig:architecture}
\end{figure}
